\chapter{Preparation}
\label{chapter:preparation}

\begin{readerguide}
 This section is absolutely key before you start an assessment! 
\end{readerguide}


In this chapter, we discuss various steps that we want to do \emph{before} we dive into assessing a code's performance.
It summarises activities that ensure that we are well-prepared to run assessments that yield meaningful outcomes. 
We typical refer to this stage as a `pre-assessment', i.e., we work through a series of steps to document our work from choosing/receiving a characteristic benchmark of a codebase, through to verifying that we can run the benchmark. 
We see this as a necessary barrier of entry because a code submitted to an assessment service should be documented sufficiently such that the analyst can initially treat the code as a black box, and get setup almost `without thinking'.


There is a checklist in Appendix~\ref{appendix:submission-checklist} which summarises some of the key information you might want to ask from a code before you start any benchmarking or analysis.
In line with the present section, this checklist is high level, but may provide an idea of the minimal information an assessor requires about a code before we get started.
It is almost a customer-facing summary of this chapter's content. 
